\documentclass[10pt,a4paper]{amsart}
\usepackage[margin=19mm]{geometry}
\usepackage{amsmath,amssymb,mathtools}
\usepackage[scr=rsfs,cal=euler]{mathalfa}
\usepackage{xfrac}
\usepackage[unicode,hidelinks,bookmarks=false]{hyperref}
\newcommand{\ie}{i.e.\ }
\newcommand{\cf}{cf.\ }
\newcommand{\resp}{respectively\ }
\newcommand{\Subsection}{\subsection}
\providecommand{\nobreakdash}{\nobreak\hbox{-}}
\setlength{\emergencystretch}{2em}
\multlinegap0pt
\usepackage[shortlabels]{enumitem}
\usepackage{amssymb}
\usepackage{dsfont}
\usepackage{xcolor}
\usepackage{braket}
\usepackage{algorithm}\usepackage{algpseudocode}
\newtheorem{lemma}{Lemma}
\newcommand{\bfB}{\mathbf{B}}
\newcommand{\bfC}{\mathbf{C}}
\newcommand{\bfF}{\mathbf{F}}
\newcommand{\bfH}{\mathbf{H}}
\newcommand{\bfI}{\mathbf{I}}
\newcommand{\bfP}{\mathbf{P}}
\newcommand{\bfR}{\mathbf{R}}
\newcommand{\bfT}{\mathbf{T}}
\newcommand{\cF}{\mathcal{F}}
\newcommand{\sifat}[1]{\marginpar{+}{\bf Sifat's remark}: {\em #1}}
\title{Communication-Efficient Approximate Gradient Coding}
\author{Sifat Munim, Aditya Ramamoorthy}
\date{Source: arXiv:2603.22514v2 (CC BY 4.0)}
\begin{document}
\maketitle

\noindent\emph{Source and licence.} Communication-Efficient Approximate Gradient Coding. Version arXiv:2603.22514v2 (CC BY 4.0), distributed by arXiv under the Creative Commons Attribution 4.0 International licence, http://creativecommons.org/licenses/by/4.0/, as recorded in the arXiv licence metadata for this exact version and shown on the abstract page https://arxiv.org/abs/2603.22514v2. Full licence notice: \texttt{licenses/arxiv-2603.22514-CC-BY-4.0.txt}.\par\smallskip

\noindent\textit{Editorial scope.} Counted unit: the article's lemma lemma:Cost\_Hadamard\_Convergence (source0.tex lines 5233--5241). The article's complete original proof of it is delivered with it (source0.tex lines 5242--5595, one \texttt{proof} environment of 354 lines). No mathematics is written, repaired or re-derived: the statement and the proof are the article's own lines, in the article's own notation, and no retained line is substituted or re-flowed.  Retained uncounted as the source-local context the proof needs: lines 1890--1890.  Nothing was omitted from any retained span, so the package carries no omission record.  No retained line contains a citation, so the wrapper carries no bibliography and no bibliography entry.  No retained line contains a figure environment, an image-inclusion command or drawing code, so no image is delivered and the package declares no asset dependency.  Editorial formatting only, in addition: an AMS \texttt{amsart} preamble that restates the article's own declarations and macros used by the retained lines, namely the environments \texttt{lemma} on separate counters, together with the standard packages those lines need.  The retained environments print in this document as Lemma 1 in order of appearance, whereas the article prints the same items with the numbering of its own full document.  The complete native source of the article as distributed in the arXiv e-print for this exact version is bundled unchanged as \texttt{sources/197022/source0.tex}.
\begin{equation} \label{eq:hada_coset_cons} \mathbf{B}=\mathbf{H}_m^T\otimes\mathbf{C}. \end{equation} 

\begin{lemma}
\label{lemma:Cost_Hadamard_Convergence}

Suppose that the assignment matrix $\mathbf{A}$ is the bi-adjacency matrix of a $(k, m, \delta)$ coset bipartite graph such that $k = p^a$, where $p$ is a prime, $a \in \mathbb{Z}_{\ge 1}$ and $p \nmid \delta$. Then, for the construction as in $\eqref{eq:hada_coset_cons}$, we have 

\[\mathbb{E}_{\mathcal{F}}[\mathbf{B}_{\cF} \mathbf{R}]= \alpha_H \mathbf{F},\]

where $\alpha_H$ is a nonzero constant.
\end{lemma}

\begin{proof}

Here, 
$\mathbf{B}^T \mathbf{B} = (\mathbf{H}_m^T \otimes \mathbf{C})^T  (\mathbf{H}_m^T \otimes \mathbf{C}) = \mathbf{H}_m\mathbf{H}_m^T \otimes \mathbf{C}^T \mathbf{C} = m \mathbf{I}_m \otimes \mathbf{C}^T \bfC$. For $\cF =\emptyset$, we assume $\bfB_{\cF} \bfR = \mathbf{0}_{mk \times m}$. For non-empty $\cF$,  let $\bfP_{\cF} = \bfI_n(:, \cF)$. We have, 
\begin{align*}
\mathbf{B}_{\mathcal{F}}\mathbf{R} &= \bfB_{\cF}(\bfB_{\cF}^T\bfB_{\cF})^{-1} \bfB_{\cF}^T \mathbf{F}\\
&=\mathbf{B} \mathbf{P}_{\mathcal{F}}(\mathbf{P}_\mathcal{F}^T \mathbf{B}^T\mathbf{B}\mathbf{P}_\mathcal{F})^{-1}\mathbf{P}_{\mathcal{F}}^T \bfB^T\mathbf{F}\\
\end{align*}

For each $i\in[m]$, let
\[
\mathcal{F}_i
:=
\left\{
j\in[k] : (i-1)k+j\in\mathcal{F}
\right\}.
\]
Thus, $\mathcal{F}_i$ denotes the indices of the non-straggling workers within the $i$-th block of $k$ workers. For non-empty $\cF_i$, let $\bfP'_{\cF_i} = \bfI_k(:, \cF_i)$. Thus, when $\cF_i \ne \emptyset$ for all $i \in [m]$,

\[
\mathbf{P}_{\mathcal{F}}
=
\begin{bmatrix}
\mathbf{P}'_{\mathcal{F}_1}
    & \mathbf{0}_{k\times |\mathcal{F}_2|}
    & \cdots
    & \mathbf{0}_{k\times |\mathcal{F}_m|}\\
\mathbf{0}_{k\times |\mathcal{F}_1|}
    & \mathbf{P}'_{\mathcal{F}_2}
    & \cdots
    & \mathbf{0}_{k\times |\mathcal{F}_m|}\\
\vdots
    & \vdots
    & \ddots
    & \vdots\\
\mathbf{0}_{k\times |\mathcal{F}_1|}
    & \mathbf{0}_{k\times |\mathcal{F}_2|}
    & \cdots
    & \mathbf{P}'_{\mathcal{F}_m}
\end{bmatrix}.
\]
% \sifat{Does it handle the case when $\cF_i$ is empty?}

Let $\mathbf{G}:=\mathbf{C}^T\mathbf{C}.$ Since $\bfB^T\bfB = m \bfI_m \otimes \bfC^T \bfC$, it follows that
\begin{align*}
\mathbf{P}_{\mathcal{F}}^T\mathbf{B}^T\mathbf{B}
\mathbf{P}_{\mathcal{F}}
& = m
\begin{bmatrix}
\mathbf{P}'^{T}_{\mathcal{F}_1} \mathbf{G}\mathbf{P}'_{\mathcal{F}_1}
& & \mathbf{0}\\
&
\ddots
&\\
\mathbf{0}
& &
\mathbf{P}'^{T}_{\mathcal{F}_m}\mathbf{G}\mathbf{P}'_{\mathcal{F}_m}
\end{bmatrix}
\end{align*}
 Therefore,
\[
\mathbf{P}_{\mathcal{F}}
(\mathbf{P}_{\mathcal{F}}^T\mathbf{B}^T\mathbf{B}
\mathbf{P}_{\mathcal{F}})^{-1}
\mathbf{P}_{\mathcal{F}}^T
=
\frac{1}{m}
\begin{bmatrix}
\mathbf{P}'_{\mathcal{F}_1}
\left(
\mathbf{P}'^{T}_{\mathcal{F}_1}
\mathbf{G}
\mathbf{P}'_{\mathcal{F}_1}
\right)^{-1}
\mathbf{P}'^{T}_{\mathcal{F}_1} & & \mathbf{0}\\
& \ddots &\\
\mathbf{0} & & \mathbf{P}'_{\mathcal{F}_m}
\left(
\mathbf{P}'^{T}_{\mathcal{F}_m}
\mathbf{G}
\mathbf{P}'_{\mathcal{F}_m}
\right)^{-1}
\mathbf{P}'^{T}_{\mathcal{F}_m}
\end{bmatrix}.
\]


For $i\in[m]$, let
% \[
% \mathbf{T}_{\mathcal{F}_i}
% :=
% \mathbf{P}_{\mathcal{F}_i}
% \left(
% \mathbf{P}_{\mathcal{F}_i}^T
% \mathbf{G}
% \mathbf{P}_{\mathcal{F}_i}
% \right)^{-1}
% \mathbf{P}_{\mathcal{F}_i}^T,
% \]
% where $\mathbf{T}_{\mathcal{F}_i}:=\mathbf{0}_{k\times k}$
% when $\mathcal{F}_i=\emptyset$. 



\[
\mathbf{T}_{\mathcal{F}_i} :=
\begin{cases}
\mathbf{P}'_{\mathcal{F}_i}
\left(
\mathbf{P}'^{T}_{\mathcal{F}_i}
\mathbf{G}
\mathbf{P}'_{\mathcal{F}_i}
\right)^{-1}
\mathbf{P}'^{T}_{\mathcal{F}_i}, & \mathcal{F}_i \neq \emptyset, \\[0.5em]
\mathbf{0}_{k \times k}, & \mathcal{F}_i = \emptyset.
\end{cases}
\]


% \sifat{say that this defn work for any non-empty and also empty F}

Since for $\cF_i = \emptyset$, $\mathbf{T}_{\cF_i} = \mathbf{0}_{k \times k}$, for any $\cF \subseteq [n]$, 

\[\mathbf{B}_{\cF} \bfR = \frac{1}{m}\mathbf{B}\begin{bmatrix}
\mathbf{T}_{\cF_1} & & \mathbf{0}\\
& \ddots &\\
\mathbf{0} & & \mathbf{T}_{\cF_m} 
\end{bmatrix} \bfB^T \bfF\]


Since each worker is independently non-straggling with probability
$1-q$, the random sets $\mathcal{F}_1,\ldots,\mathcal{F}_m$ are
independent. Let
$
\overline{\mathbf{T}}
:=
\mathbb{E}_{\mathcal{F}}
[\mathbf{T}_{\mathcal{F}_i}]
.$ Thus,
\[
\mathbb{E}_{\mathcal{F}}
\left[
\bfB_{\cF}\bfR
\right]
=
\frac{1}{m} \bfB (\mathbf{I}_m\otimes\overline{\mathbf{T}}) \mathbf{B}^T \bfF.
\]

Next, we show that $\overline{\mathbf{T}}$ is circulant. Let
$\mathbf{Q}\in\{0,1\}^{k\times k}$ be the permutation matrix
corresponding to a cyclic shift by one position. Since $\mathbf{C}$
is circulant,
$
\mathbf{Q}^T\mathbf{C}\mathbf{Q}=\mathbf{Q}\mathbf{C}\mathbf{Q}^T =\mathbf{C}.
$ From this we can infer that $\mathbf{Q}^T\mathbf{G}\mathbf{Q} = \mathbf{Q}\mathbf{G}\mathbf{Q}^T
= \mathbf{G}$.


For $\mathcal{F}_i\subseteq[k]$, let $\mathcal{F}_i'$
denote the set obtained by cyclically shifting the elements of
$\mathcal{F}_i$ by one position. Then,
$\mathbf{Q}\mathbf{P}'_{\mathcal{F}_i}$ is an equivalent selection
matrix for $\mathcal{F}_i'$. 
% Hence,
% \[
% \mathbf{T}_{\mathcal{S}'}
% =
% \mathbf{Q}\mathbf{T}_{\mathcal{S}}\mathbf{Q}^T.
% \]
% Since the workers straggle independently, 
% \[
% \mathcal{S}'\overset{d}{=}\mathcal{S}.
% \]
Now,
\begin{align*}
\mathbf{Q} \bfT_{\mathcal{F}_i} \mathbf{Q}^T &= \mathbf{Q} \mathbf{P}'_{\mathcal{F}_i}
\left(
\mathbf{P}'^{T}_{\mathcal{F}_i}
\mathbf{G}
\mathbf{P}'_{\mathcal{F}_i}
\right)^{-1}
\mathbf{P}'^{T}_{\mathcal{F}_i} \mathbf{Q}^T\\
&= \mathbf{Q} \mathbf{P}'_{\mathcal{F}_i}
\left(
\mathbf{P}'^{T}_{\mathcal{F}_i}\mathbf{Q}^T\mathbf{Q}
\mathbf{G}\mathbf{Q}^T\mathbf{Q}
\mathbf{P}'_{\mathcal{F}_i}
\right)^{-1}
\mathbf{P}'^{T}_{\mathcal{F}_i} \mathbf{Q}^T\\
&= \mathbf{Q} \mathbf{P}'_{\mathcal{F}_i}
\left(
\mathbf{P}'^{T}_{\mathcal{F}_i}\mathbf{Q}^T
\mathbf{G}\mathbf{Q}
\mathbf{P}'_{\mathcal{F}_i}
\right)^{-1}
\mathbf{P}'^{T}_{\mathcal{F}_i} \mathbf{Q}^T\\
& = \bfT_{\cF_i'}.
\end{align*}
% \begin{align*}

% \end{align*}
% \[
% \mathbf{Q}\mathbf{T}_{\mathcal{S}}\mathbf{Q}^T
% \overset{d}{=}
% \mathbf{T}_{\mathcal{S}}.
% \]

Since $\cF_i  \overset{d}{=} \cF_i'$, we have $\bfT_{\cF_i} \overset{d} {=} \bfT_{\cF_i'}$.   Thus, $\mathbb{E}_{\cF}\left[\bfT_{\cF_i}\right] = \mathbb{E}_{\cF}\left[\bfT_{\cF_i'}\right]$. Consequently, $
\mathbf{Q}\overline{\mathbf{T}}\mathbf{Q}^T
=
\overline{\mathbf{T}}.
$
So, by noting that $\mathbf{Q}$ corresponds to a cyclic shift by one position, we can conclude that $\overline{\mathbf{T}}$ is circulant. So,
\begin{align*}
\mathbb{E}_{\mathcal{F}}
[\mathbf{B}_{\mathcal{F}}\mathbf{R}]
&=
\frac{1}{m} \bfB (\mathbf{I}_m\otimes\overline{\mathbf{T}}) \mathbf{B}^T \bfF\\
&=
\frac{1}{m}(\mathbf{H}_m^T\otimes\mathbf{C})
\left(\mathbf{I}_m\otimes\overline{\mathbf{T}}
\right)
(\mathbf{H}_m\otimes\mathbf{C}^T)(\bfI_m \otimes \mathbf{1}_k)\\
&=\frac{1}{m} \bfH_m^T\bfH_m \otimes \bfC \overline{\mathbf{T}} \bfC^T\mathbf{1}_k = \bfI_m \otimes  \bfC \overline{\mathbf{T}} \bfC^T\mathbf{1}_k .
\end{align*}

Let
$
\tilde{\mathbf{T}}:=
\mathbf{C}\overline{\mathbf{T}}\mathbf{C}^T.$
Since $\mathbf{C}$ and $\overline{\mathbf{T}}$ are circulant,
$\tilde{\mathbf{T}}$ is also circulant. Therefore, $\mathbf{1}_k$ is an
eigenvector of $\tilde{\mathbf{T}}$, and there exists a constant $\alpha_H$
such that $\tilde{\mathbf{T}}\mathbf{1}_k=\alpha_H\mathbf{1}_k.$


It remains to show that $\alpha_H$ is nonzero. Since $\mathbf{C}$ is
invertible, $\mathbf{G}=\mathbf{C}^T\mathbf{C}$ is positive definite.
Thus, for any non-empty $\mathcal{F}_i$,
$
\mathbf{1}_k^T\mathbf{T}_{\mathcal{F}_i}\mathbf{1}_k
=
\mathbf{1}_{|\mathcal{F}_i|}^T
\left(
\mathbf{P}'^{T}_{\mathcal{F}_i}
\mathbf{G}
\mathbf{P}'_{\mathcal{F}_i}
\right)^{-1}
\mathbf{1}_{|\mathcal{F}_i|}
>0.
$
Since $q<1$, $\mathcal{F}_i$ is non-empty with positive probability.
Therefore,
$
\mathbf{1}_k^T\overline{\mathbf{T}}\mathbf{1}_k>0
.$ Moreover, since $\mathbf{C}$ is $\delta$-regular,
$\mathbf{C}\mathbf{1}_k=\mathbf{C}^T\mathbf{1}_k
=\delta\mathbf{1}_k$. Hence,
$
\alpha_H k
=
\mathbf{1}_k^T\widetilde{\mathbf{T}}\mathbf{1}_k
=
\delta^2
\mathbf{1}_k^T\overline{\mathbf{T}}\mathbf{1}_k
>0.
$ Consequently, $\alpha_H$ is nonzero.






% Recall that
% \[
% \mathbf{f}_i
% =
% \mathbf{e}_i\otimes\mathbf{1}_k,
% \qquad i\in[m].
% \]
% Thus,
% \[
% (\mathbf{I}_m\otimes\mathbf{M})\mathbf{f}_i
% =
% \mathbf{e}_i\otimes
% \mathbf{M}\mathbf{1}_k
% =
% \alpha\mathbf{f}_i.
% \]
% Since
% $\mathbf{F}=[\mathbf{f}_1|\cdots|\mathbf{f}_m]$, it follows that
% \[
% \mathbb{E}_{\mathcal{F}}
% [\mathbf{B}_{\mathcal{F}}\mathbf{R}]
% =
% \alpha\mathbf{F}.
% \]

% Finally, $\alpha$ is non-zero. Since $\mathbf{C}$ is invertible,
% $\mathbf{G}=\mathbf{C}^T\mathbf{C}$ is positive definite, and thus,
% for every non-empty $\mathcal{S}$,
% \[
% \mathbf{1}_k^T\mathbf{T}_{\mathcal{S}}\mathbf{1}_k
% =
% \mathbf{1}_{|\mathcal{S}|}^T
% \left(
% \mathbf{P}_{\mathcal{S}}^T
% \mathbf{G}
% \mathbf{P}_{\mathcal{S}}
% \right)^{-1}
% \mathbf{1}_{|\mathcal{S}|}
% >0.
% \]
% Since $q<1$, a non-empty set of non-stragglers occurs with positive
% probability. Hence,
% \[
% \mathbf{1}_k^T\overline{\mathbf{T}}\mathbf{1}_k>0.
% \]
% Also, $\mathbf{C}\mathbf{1}_k=\delta\mathbf{1}_k$ and
% $\mathbf{C}^T\mathbf{1}_k=\delta\mathbf{1}_k$. Consequently,
% \[
% \mathbf{1}_k^T\mathbf{M}\mathbf{1}_k
% =
% \delta^2
% \mathbf{1}_k^T\overline{\mathbf{T}}\mathbf{1}_k
% >0,
% \]
% which implies $\alpha>0$. Therefore,
% \[
% \mathbb{E}_{\mathcal{F}}
% [\mathbf{B}_{\mathcal{F}}\mathbf{R}]
% =
% \alpha\mathbf{F},
% \]
% where $\alpha$ is a non-zero constant.


















\end{proof}

\end{document}
